% Bewertungspaket zum Gesamttest «Statik» — Quelle des PDFs (python3 scripts/build-lp-pdf.py statik).
% Ganze Punkte je Teilschritt; Werte und Folgefehler-Fälle mit python3 nachgerechnet (05.10.2026).
% Aufbau, Auftrag an die KI und Punktarten wie im Leitprogramm Energie, dazu (S) für die Kräfteskizze.
\documentclass[11pt]{article}
\usepackage{../lp-druck}
\renewcommand{\lpname}{Statik}
\renewcommand{\lpteilgebiet}{4.4}
\renewcommand{\lpfuss}{Bewertungspaket}
\newcolumntype{P}{>{\centering\arraybackslash}p{10mm}}
\newenvironment{raster}{\par\smallskip\noindent\tabularx{\linewidth}{@{}>{\raggedright\arraybackslash}p{38mm}P X@{}}\toprule
  \small\textsc{Teilschritt} & \small\textsc{P} & \small\textsc{Musterlösung}\\\midrule}{\bottomrule\endtabularx\par}
\newcommand{\fehler}[1]{\par\smallskip{\small\textcolor{lprot}{\bfseries Typische Fehler:} #1}}
\newcommand{\E}{\,{\footnotesize\bfseries\color{lpgruen}(E)}}
\newcommand{\B}{\,{\footnotesize\bfseries(B)}}
\newcommand{\W}{\,{\footnotesize\bfseries(W)}}
\newcommand{\Sk}{\,{\footnotesize\bfseries(S)}}
\newcommand{\A}{\,{\footnotesize\bfseries\color{lpgruen}(A)}}
\begin{document}
\lpkopf{Bewertungspaket zum Gesamttest}{}

\begin{lpachtung}{Erst lösen, dann öffnen}
Dieses Paket enthält die Musterlösung. Wer es vor dem Lösen liest, misst am Ende nur, wie gut das Abschreiben gelingt.
\end{lpachtung}

\section*{1 · So gehst du vor}
\begin{enumerate}[leftmargin=*,itemsep=2pt]
\item \textbf{Gesamttest lösen} — vollständig, mit Rechenweg, ohne dieses Paket.
\item \textbf{Fotografieren oder scannen}, gut lesbar, eine Seite pro Bild. \textbf{Kein Name} auf den Bildern.
  Handyfotos tragen oft unsichtbar Standort, Zeit und Gerät mit (Metadaten). Lade darum lieber einen Scan oder ein
  Bildschirmfoto des Fotos hoch, oder entferne vorher die Standortangaben.
\item \textbf{Bewerten lassen.} Ob du dafür eine KI nutzt und welche, entscheidest du selbst und in eigener Verantwortung —
  je nachdem, was dir aufgrund deines Alters und deiner persönlichen Situation erlaubt ist (Altersgrenzen und
  Nutzungsbedingungen des Anbieters, allenfalls Einverständnis deiner Eltern). Lade nur deine Lösung und dieses PDF hoch,
  keine persönlichen Daten, und füge den Auftrag aus Abschnitt~2 ein. Ohne KI kannst du dich mit Abschnitt~3 selbst bewerten.
\item \textbf{Rückmeldung prüfen.} Stimmt, was die KI aus deiner Handschrift gelesen hat? Eine KI kann sich irren —
  im Zweifel gilt das Raster in Abschnitt~3.
\item \textbf{Gezielt wiederholen} nach Abschnitt~4.
\end{enumerate}

\needspace{14\baselineskip}
\section*{2 · Auftrag an die KI}
\begin{tcolorbox}[colback=lplinie!25,colframe=lplinie,boxrule=0.5pt,arc=1mm,fontupper=\small]
Du bist Korrektorin oder Korrektor für Physik an einer Schweizer Berufsmaturitätsschule. Im Anhang findest du (1)~das Bewertungspaket zum Gesamttest «Statik» mit Musterlösung und Punkteraster und (2)~Fotos meiner handschriftlichen Lösung.\par\medskip
Geh so vor:\par
1. Schreib zu jeder Aufgabe G1 bis G6 zuerst ab, was du in meiner Lösung liest — Rechenweg und Ergebnis. Was du nicht sicher lesen kannst, markierst du mit [unsicher]. Nicht raten.\par
2. Vergib die Punkte genau nach dem Raster im Bewertungspaket (Abschnitt 3), ganze Punkte je Zeile. Ziehe einen Fehler nur einmal ab: Rechne ich mit einem falschen Zwischenergebnis richtig weiter, gibt es die Folgepunkte — auch wenn das Folgeergebnis unplausibel ist; eine fehlende Plausibilitätsprobe kostet keinen zusätzlichen Punkt. Die Zeilen «Typische Fehler» sagen, welche Rasterzeile bei diesem Fehler 0 Punkte gibt — sie sind kein zusätzlicher Abzug.\par
3. Andere richtige Lösungswege zählen voll. Gleichwertige Schreibweisen zählen voll: Dezimalpunkt oder Dezimalkomma, andere Einheit oder Vorsilbe ($0.12\;\text{m} = 12\;\text{cm}$; $\text{m/s}$ statt $\text{km/h}$, wenn richtig umgerechnet), Zehnerpotenz oder ausgeschriebene Zahl, sinnvoll gerundete Werte (drei signifikante Stellen). Zeilen mit (E) sind Ergebnispunkte: Sie verlangen das Ergebnis \emph{mit richtiger Einheit}, auch ohne Weg; ohne oder mit falscher Einheit gibt es diesen Punkt nicht. Zeilen mit (W) gibt es nur mit nachvollziehbarem Weg (Formel, dann Werte mit Einheiten); ein Zahlenwert darin braucht ebenfalls die Einheit. Zeilen mit (B) sind Begründungen: Es zählt die fachliche Aussage in eigenen Worten, eine Formel ist nicht nötig; andere richtige Formulierungen zählen voll. Zeilen mit (S) sind Skizzenpunkte: Es zählt, dass die verlangten Kräfte als Pfeile mit Namen und richtiger Richtung eingezeichnet sind; Länge und Schönheit zählen nicht. Werte aus sinnvoll gerundeten Zwischenergebnissen zählen, wenn sie höchstens rund 2\,\% vom Raster abweichen.\par
4. Begründe jeden Abzug in einem Satz und nenne die Fehlerart (zum Beispiel «Sinus statt Kosinus»). Schreib die Musterlösung nicht ab — zeig mir, wo mein Fehler liegt.\par
5. Gib zum Schluss eine Tabelle «Aufgabe | Punkte | Begründung», die Gesamtpunktzahl (von 25) und — nach der Selbsteinschätzung in Abschnitt 4 — welches Kapitel des Leitprogramms ich wiederholen soll. Keine Note.\par
6. Fehlt ein Foto oder ist eine Aufgabe unlesbar, sag es, statt sie zu bewerten.
\end{tcolorbox}

\needspace{16\baselineskip}
\section*{3 · Musterlösung und Punkteraster}
\textbf{Allgemein:} Ganze Punkte je Zeile. Ein Fehler wird nur einmal abgezogen: Wer mit einem falschen Zwischenergebnis
richtig weiterrechnet, bekommt die folgenden Zeilen (Folgefehler) — auch wenn das Folgeergebnis unplausibel ist; eine fehlende
Plausibilitätsprobe kostet keinen zusätzlichen Punkt. Gleichwertige Wege und Schreibweisen zählen voll.
In der Spalte P: \textbf{(E)}~= Ergebnispunkt — es zählt das Ergebnis \emph{mit richtiger Einheit}, auch ohne Weg; ohne oder mit
falscher Einheit gibt es ihn nicht. \textbf{(W)}~= Wegpunkt — nur mit nachvollziehbarem Weg (Formel, dann Werte mit Einheiten);
steht in einer Weg-Zeile ein Zahlenwert, braucht auch er die Einheit. \textbf{(B)}~= Begründungspunkt — es zählt die fachliche
Aussage in eigenen Worten, eine Formel ist nicht nötig. \textbf{(S)}~= Skizzenpunkt — die verlangten Kräfte als Pfeile mit Namen
und richtiger Richtung. Folgewerte aus gerundeten Zwischenergebnissen zählen, wenn sie höchstens rund 2\,\% abweichen.
«Typische Fehler» nennt die Zeilen mit 0 Punkten und die Punktzahl, die bleibt.

\begin{lpaufgabe}{G1}{Kraft als Vektor}{4 P}
\begin{raster}
a) Kraft & 1\B & Eine Kraft verformt einen Körper oder ändert seinen Bewegungszustand; festgelegt durch Betrag, Richtung und Angriffspunkt (beides nötig)\\
b) Komponenten & 1\E & \(F_x = F \cdot \cos\varphi = 5.0\;\text{kN} \cdot \cos 220^\circ \approx -3.83\;\text{kN}\), \(F_y = F \cdot \sin\varphi \approx -3.21\;\text{kN}\) (beide mit Minus; gleichwertig \(-3830\;\text{N}\), \(-3210\;\text{N}\))\\
c) \(F_y\) & 1\E & \(F_y = \sqrt{F^2 - F_x^2} = \sqrt{(2.6\;\text{kN})^2 - (1.0\;\text{kN})^2} = 2.4\;\text{kN}\) (auch \(\pm 2.4\;\text{kN}\))\\
d) zwei Richtungen & 1\B & \(F_y = +2.4\;\text{kN}\) oder \(-2.4\;\text{kN}\): nach rechts oben, \(\varphi \approx 67.4^\circ\), oder nach rechts unten, \(\varphi \approx 292.6^\circ\) (gleichwertig \(-67.4^\circ\)); im Quadrat \(F_y^2\) fällt das Vorzeichen weg, darum passen beide (beide Winkel und Grund nötig)\\
\end{raster}
\fehler{b) Sinus und Kosinus vertauscht (\(-3.21\;\text{kN}\) und \(-3.83\;\text{kN}\)) oder ohne Minus: b) 0 \(\to\) 3 von 4\,P. Taschenrechner im Bogenmass (\(4.98\;\text{kN}\), \(0.44\;\text{kN}\)): b) 0 \(\to\) 3 von 4\,P.
c) \(2.6\;\text{kN} - 1.0\;\text{kN} = 1.6\;\text{kN}\) (subtrahiert): c) 0; d) folgerichtig mit \(\arctan 1.6\) zählt \(\to\) 3 von 4\,P.}
\end{lpaufgabe}

\begin{lpaufgabe}{G2}{Drei Kräfte am Ring}{4 P}
\begin{raster}
Ansatz & 1\W & komponentenweise addieren: \(F_{\text{res},x} = \sum F \cdot \cos\varphi\), \(F_{\text{res},y} = \sum F \cdot \sin\varphi\)\\
a) Komponenten & 1\E & \(F_{\text{res},x} = 300\;\text{N} + 250\;\text{N} \cdot \cos 210^\circ \approx 83.5\;\text{N}\), \(F_{\text{res},y} = 200\;\text{N} + 250\;\text{N} \cdot \sin 210^\circ = 75\;\text{N}\)\\
b) Resultierende & 1\E & \(F_\text{res} = \sqrt{(83.5\;\text{N})^2 + (75\;\text{N})^2} \approx 112\;\text{N}\), Richtung \(\varphi = \arctan\dfrac{75}{83.5} \approx 41.9^\circ\) (beides nötig)\\
c) Gleichgewicht & 1\B & gleich gross und entgegengesetzt: \(112\;\text{N}\) unter \(\varphi \approx 221.9^\circ\) (gleichwertig: nach links unten, \(41.9^\circ\) unter der negativen \(x\)-Achse), weil dann \(\vec{F}_\text{res} = \vec{0}\) ist\\
\end{raster}
\fehler{Beträge addiert (\(750\;\text{N}\)): Ansatz, a) und b) 0; c) folgerichtig «entgegengesetzt» zählt \(\to\) 1 von 4\,P.
\(F_3\) ohne Vorzeichen (\(516.5\;\text{N}\), \(325\;\text{N}\)): a) 0; b) folgerichtig \(\approx 610\;\text{N}\), \(32.2^\circ\) und c) folgerichtig zählen \(\to\) 3 von 4\,P.}
\end{lpaufgabe}

\begin{lpaufgabe}{G3}{Lenkrad, \(r = 0.19\;\text{m}\)}{4 P}
\begin{raster}
a) senkrecht & 1\E & \(M = F \cdot r = 30\;\text{N} \cdot 0.19\;\text{m} = 5.7\;\text{Nm}\)\\
b) schräg & 1\E & \(M = F \cdot l \cdot \sin\alpha = 30\;\text{N} \cdot 0.19\;\text{m} \cdot \sin 50^\circ \approx 4.37\;\text{Nm}\)\\
c) kleiner & 1\E & \(F = \dfrac{M}{r} = \dfrac{5.7\;\text{Nm}}{0.15\;\text{m}} = 38\;\text{N}\)\\
d) Lastwagen & 1\B & Grosser Radius heisst grosser Hebelarm: Für das nötige Drehmoment genügt nach \(F = \dfrac{M}{r}\) eine kleinere Kraft\\
\end{raster}
\fehler{b) Kosinus statt Sinus (\(3.66\;\text{Nm}\)): b) 0 \(\to\) 3 von 4\,P. c) \(5.7\;\text{Nm} \cdot 0.15\;\text{m}\) (mal statt durch): c) 0 \(\to\) 3 von 4\,P.}
\end{lpaufgabe}

\begin{lpaufgabe}{G4}{Auto, \(1200\;\text{kg}\), \(\mu_H = 0.35\)}{5 P}
\begin{raster}
a) Skizze & 1\Sk & Gewichtskraft senkrecht nach unten (im Schwerpunkt), Normalkraft senkrecht zur Strasse weg von ihr, Haftreibung längs der Strasse hangaufwärts; alle drei mit Namen. Eine eingezeichnete Zerlegung der Gewichtskraft schadet nicht\\
b) Neigung & 1\E & Grenze \(F_H = \mu_H \cdot F_N\), also \(\tan\alpha = \mu_H\): \(\alpha = \arctan 0.35 \approx 19.3^\circ\)\\
c) Normalkraft & 1\E & \(F_N = m \cdot g \cdot \cos\alpha = 1200\;\text{kg} \cdot 9.81\;\text{m/s}^2 \cdot \cos 19.3^\circ \approx 11.1\;\text{kN}\)\\
c) Haftreibung & 1\E & an der Grenze \(F_R = \mu_H \cdot F_N \approx 3.89\;\text{kN}\) (gleichwertig \(F_R = F_H = m \cdot g \cdot \sin\alpha\))\\
d) beladen & 1\B & bei derselben Neigung: Hangabtrieb und grösste Haftreibung wachsen beide mit der Masse, die Masse kürzt sich, \(\tan\alpha = \mu_H\)\\
\end{raster}
\fehler{a) Normalkraft senkrecht nach oben statt senkrecht zur Strasse, oder Reibung hangabwärts: a) 0 \(\to\) 4 von 5\,P.
b) \(\arcsin 0.35 \approx 20.5^\circ\): b) 0; c) folgerichtig (\(\approx 11.0\;\text{kN}\), \(\approx 3.86\;\text{kN}\)) zählt \(\to\) 4 von 5\,P.
c) \(F_N = F_G \approx 11.8\;\text{kN}\) oder \(F_R = \mu_H \cdot F_G \approx 4.12\;\text{kN}\): die betroffene Zeile 0 \(\to\) 4 von 5\,P.}
\end{lpaufgabe}

\begin{lpaufgabe}{G5}{Malerbrett, \(3\;\text{m}\)}{5 P}
\begin{raster}
Ansatz & 1\W & Momente um A (\(F_A\) fällt heraus): \(F_B \cdot 3\;\text{m} = 100\;\text{N} \cdot 0.5\;\text{m} + 150\;\text{N} \cdot 1.5\;\text{m} + 700\;\text{N} \cdot 2\;\text{m}\)\\
a) \(F_B\) & 1\E & \(F_B = \dfrac{1675\;\text{Nm}}{3\;\text{m}} \approx 558\;\text{N}\)\\
a) \(F_A\) & 1\E & \(F_A = 100\;\text{N} + 150\;\text{N} + 700\;\text{N} - F_B \approx 392\;\text{N}\) (gleichwertig: Momente um B)\\
b) Drehachse & 1\B & Die Auflagerkraft dort hat keinen Hebelarm und fällt aus der Momentengleichung heraus; es bleibt eine Gleichung mit einer Unbekannten\\
c) wie weit & 1\E & \(F_B = 750\;\text{N}\): \(750\;\text{N} \cdot 3\;\text{m} = 50\;\text{Nm} + 225\;\text{Nm} + 700\;\text{N} \cdot x\), \(x \approx 2.82\;\text{m}\) von A\\
\end{raster}
\fehler{Eigengewicht des Bretts vergessen (\(F_B \approx 483\;\text{N}\), \(F_A \approx 317\;\text{N}\)): Ansatz und \(F_B\) 0; \(F_A\) und c) folgerichtig (\(x \approx 3.14\;\text{m}\), also bis ans Ende) zählen \(\to\) 3 von 5\,P.
Eigengewicht bei B statt in der Mitte (\(F_B \approx 633\;\text{N}\)): Ansatz 0; Folgewerte zählen \(\to\) 4 von 5\,P.}
\end{lpaufgabe}

\begin{lpaufgabe}{G6}{Gleichgewicht}{3 P}
\begin{raster}
a) Bedingungen & 1\B & Er bleibt in Ruhe: Die Summe aller Kräfte ist null (\(\sum \vec{F} = \vec{0}\)) \emph{und} die Summe aller Drehmomente ist null (\(\sum M = 0\)) — beide nötig\\
b) Stab & 1\B & Nein: Die Kräfte heben sich auf (\(\sum \vec{F} = \vec{0}\), er verschiebt sich nicht), aber sie greifen an verschiedenen Enden an und drehen beide gleich herum: \(\sum M \ne 0\), der Stab beginnt sich zu drehen\\
c) Wippe & 1\E & \(m_1 \cdot r_1 = m_2 \cdot r_2\): \(r_2 = \dfrac{25\;\text{kg} \cdot 1.6\;\text{m}}{40\;\text{kg}} = 1.0\;\text{m}\)\\
\end{raster}
\fehler{b) «Ja, die Kräfte heben sich auf» ohne Momente: b) 0 \(\to\) 2 von 3\,P. c) umgekehrt (\(2.56\;\text{m}\)): c) 0 \(\to\) 2 von 3\,P.}
\end{lpaufgabe}

\needspace{14\baselineskip}
\section*{4 · Selbsteinschätzung}
\begin{tabularx}{\linewidth}{@{}l X@{}}\toprule
22 – 25 P & Die geprüften Teile sitzen. Wo du Punkte verloren hast: das Kapitel dieser Aufgabe nochmals (Zuordnung unten).\\
17 – 21 P & Den schwächsten Teil nochmals: Simulation und Übungen des Kapitels, in dem du die meisten Punkte verloren hast.\\
11 – 16 P & Zurück zu den Kapiteln aller Aufgaben, in denen du Punkte verloren hast.\\
0 – 10 P & Zurück zu Kapitel 1 und von dort der Reihe nach weiter.\\\midrule
\multicolumn{2}{@{}p{\linewidth}@{}}{\small\textbf{Aufgabe $\to$ Kapitel} (gilt für alle Bereiche): G1 $\to$ Kapitel 1 (K1) · G2 $\to$ Kapitel 2 (K4) · G3 $\to$ Kapitel 4 (K2) · G4 $\to$ Kapitel 3 (K3, K5) · G5 $\to$ Kapitel 6 (K5) · G6 $\to$ Kapitel 5, 6 (K5)}\\\bottomrule
\end{tabularx}
\end{document}
